% 星座（综述）
% license CCBYSA3
% type Wiki

本文根据 CC-BY-SA 协议转载翻译自维基百科\href{https://en.wikipedia.org/wiki/Constellation}{相关文章}。

\begin{figure}[ht]
\centering
\includegraphics[width=6cm]{./figures/a943f80cad7e3c1f.png}
\caption{} \label{fig_Coti_1}
\end{figure}
星座是天球上的一个区域，在该区域内，一组可见的恒星形成了人们所感知的图案或轮廓，通常代表一种动物、神话题材，或无生命的物体。$^{[1]}$

最早的星座很可能是在史前时期被定义的。人们利用星座来讲述关于其信仰、经历、创世观念以及神话的故事。不同的文化和国家创造了各自的星座体系，其中一些一直延续到20世纪初，直到当今的星座体系在国际上获得统一认可。随着时间的推移，对星座的认知发生了显著变化。许多星座的大小或形状发生了改变。有些一度流行，后来却逐渐被遗忘。有些则仅局限于单一文化或国家。对星座进行命名也帮助天文学家和航海者更容易地识别恒星。$^{[2]}$

十二个（或十三个）古代星座属于黄道星座，它们横跨黄道——太阳、月球和行星运行所经过的路径。黄道的起源在历史上仍不确定；其占星学上的划分大约在公元前400年d  LKJHGFDSAjhgfdsa PODA于巴比伦或迦勒底天文学中变得突出。$^{[3]}$ 星座通过希腊文化进入西方世界，并出现在赫西俄德、欧多克索斯和阿拉托斯的著作中。由黄道星座及另外36个星座组成的传统48星座（由于阿尔戈船座被划分为三个星座，如今为38个），由出生于埃及亚历山大的希腊—罗马天文学家托勒密在其著作《天文学大成》中列出。星座的形成是大量神话的主题，其中最著名的当属拉丁诗人奥维德的《变形记》。15世纪至18世纪中叶，随着欧洲探险者开始前往南半球，远南天区的星座被陆续加入。由于罗马和欧洲文化的传播，每个星座都拥有一个拉丁名称。

1922年，国际天文学联合会正式接受了现代的88个星座列表，并于1928年采用了官方的星座边界，这些边界共同覆盖了整个天球。$^{[4][5]}$ 在天球坐标系统中，任意一个给定点都位于某一个现代星座之内。一些天文学命名体系在命名天体时会包含其所在的星座，以传达该天体在天空中的大致位置。例如，恒星的弗兰斯蒂德命名法由一个数字和星座名称的所有格形式组成。

其他被称为星群的恒星图案或组合，在严格定义下并不属于星座，但同样被观测者用来辨认夜空。星群可能是某一星座中的几颗恒星，也可能与多个星座共享恒星。星群的例子包括人马座中的茶壶，以及大熊座中的北斗七星。$^{[6][7]}$
\subsection{术语}  
“星座”一词源自晚期拉丁语 cōnstellātiō，可译为“恒星的集合”；该词于14世纪进入中古英语。$^{[8]}$ 古希腊语中表示星座的词是 ἄστρον（astron）。这些术语在历史上泛指任何可被识别的恒星图案，其外观通常与神话人物或生物、地面动物或某种物体相关联。$^{[1]}$ 随着时间推移，在欧洲天文学家之间，星座逐渐被明确界定并获得广泛认可。20世纪，国际天文学联合会（IAU）正式确认了88个星座。$^{[9]}$

当从地球上某一特定纬度观测时，若某个星座或恒星永远不会落到地平线以下，则称其为拱极星座。从北极或南极观测，位于天赤道以南或以北的所有星座都是拱极的。根据不同定义，赤道星座可以指位于北纬45°至南纬45°之间的星座$^{[10]}$，也可以指穿过黄道（或黄道带）赤纬范围、即北纬23.5°至南纬23.5°之间的星座。$^{[11][12]}$

星座中的恒星在天空中看似彼此接近，但它们通常位于距离地球各不相同的位置。由于每颗恒星都具有各自独立的运动，所有星座都会随着时间缓慢发生变化。经过数万至数十万年，熟悉的星座轮廓将变得难以辨认。$^{[13]}$天文学家可以通过精确测量单个恒星的共同自行$^{[14]}$（依赖高精度天体测量学$^{[15][16]}$）以及通过天文光谱学测定其径向速度$^{[17]}$，来预测星座在过去或未来的形态。

国际天文学联合会所承认的88个星座，以及历史上各文化所使用的星座，都是基于可观测天空中恒星分布而想象出的形象与形状。许多官方承认的星座源自古代近东和地中海地区的神话想象。$^{[18][19][page needed]}$ 其中一些故事似乎与星座的外观有关，例如猎户座被天蝎座刺杀的传说，其对应的星座在一年中的不同季节出现在天空中。$^{[20]}$
\subsection{观测}
\begin{figure}[ht]
\centering
\includegraphics[width=8cm]{./figures/479bf5216e7afbe4.png}
\caption{显示星座、银河以及黄道的星图} \label{fig_Coti_2}
\end{figure}
由于地球围绕太阳运行的轨道上，夜晚发生的时间对应于轨道上的位置不断变化，星座在一年之中其位置也会随之改变。随着地球自转方向指向东方，天球看起来仿佛向西旋转，恒星在北极星周围呈逆时针环绕，在南极星周围则呈顺时针环绕。$^{[21]}$

由于地球自转轴具有23.5°的倾角，黄道沿着黄道面在南、北半球中均匀分布，近似构成一个大圆。北天的黄道星座包括双鱼座、白羊座、金牛座、双子座、巨蟹座和狮子座；南天的黄道星座包括室女座、天秤座、天蝎座、人马座、摩羯座和宝瓶座。$^{[22][a]}$ 根据一年中的不同时间，从北纬23.5°到南纬23.5°的纬度范围内，黄道会出现在头顶正上方。夏季时，黄道在白天显得较高、夜晚较低；而在冬季，情况则正好相反，这一现象在南、北半球均一致。

由于太阳系相对于银河系具有约60°的倾角，银河系的银道面相对于黄道倾斜约60°，位于北天的金牛座与双子座之间，以及南天的天蝎座与人马座之间，银河系中心也位于这一南天区域附近。$^{[22][23]}$ 银河在天空中穿过天鹰座（靠近天赤道），以及北天的天鹅座、仙后座、英仙座、御夫座和猎户座（靠近参宿四），同时也穿过麒麟座（靠近天赤道），以及南天的船尾座、船帆座、船底座、南十字座、半人马座、南三角座和天坛座。$^{[22]}$
\subsubsection{北半球}  
北极星作为北极星，是北天球的近似中心。它属于小熊座，位于小北斗勺柄的末端。$^{[22]}$

在大约北纬35°的地区，一月时，大熊座（包含北斗七星）出现在东北方，而仙后座位于西北方。西方可见双鱼座（位于地平线上方）和白羊座。西南方的鲸鱼座接近地平线。高空偏南的位置可见猎户座和金牛座。东南方向地平线上方是大犬座。在猎户座的上方偏东方向可以看到双子座；东方方向（并逐渐接近地平线）还有巨蟹座和狮子座。除金牛座外，英仙座和御夫座也出现在头顶附近。$^{[22]}$

在相同纬度下，七月时，仙后座（位于低空）和仙王座出现在东北方。大熊座此时位于西北方。牧夫座高悬于西方天空。室女座位于西方，天秤座在西南方，天蝎座在南方。人马座和摩羯座位于东南方。天鹅座（包含北十字星）位于东方。武仙座与北冕座一同高悬于天空中。$^{[22]}$
\begin{figure}[ht]
\centering
\includegraphics[width=6cm]{./figures/58cdd9b51e42621a.png}
\caption{南十字座中的南十字星群以及半人马座中的“南天指针星”可用于寻找南天极星——南极座$\sigma$星。} \label{fig_Coti_3}
\end{figure}
\subsubsection{南半球}  
一月可见的星座包括绘架座和网罟座，分别位于水蛇座和山案座附近。$^{[24]}$

七月时，可以看到天坛座（毗邻南三角座）以及天蝎座。$^{[25]}$

靠近极星的星座包括蝘蜓座、天燕座和南三角座（靠近半人马座），以及孔雀座、水蛇座和山案座。

南极座σ星是最接近南天极的恒星，但在夜空中十分暗弱。因此，可以利用南十字座以及半人马座中的α星和β星（相对于南十字座约逆时针30°）来进行三角定位以确定南天极的位置。$^{[22]}$
\subsection{早期星座的历史}
\subsubsection{法国南部拉斯科洞穴}  
有人提出，位于法国南部拉斯科的、距今约17000年的洞穴壁画描绘了金牛座、猎户座腰带以及昴宿星团等星座。然而，这一观点并未在科学界得到普遍认可。$^{[26][27]}$
\subsubsection{美索不达米亚} 
来自美索不达米亚（今伊拉克地区）的刻石和泥板文书，其年代可追溯至公元前3000年，被认为是人类识别星座的最早、且普遍被接受的证据。$^{[28]}$ 看来，大多数美索不达米亚星座是在相对较短的时间内形成的，大约集中于公元前1300年至公元前1000年之间。美索不达米亚的星座后来出现在许多古典希腊星座体系之中。$^{[29]}$
\subsubsection{古代近东}
\begin{figure}[ht]
\centering
\includegraphics[width=6cm]{./figures/71a76fd0377888ed.png}
\caption{记录公元前164年哈雷彗星的巴比伦泥板} \label{fig_Coti_4}
\end{figure}
最古老的巴比伦恒星与星座目录可追溯至中青铜时代初期，其中尤以《每组三星》文本以及《MUL.APIN》最为著名；后者是在约公元前1000年基于更为精确的观测，对前者加以扩展和修订而成。然而，这些目录中大量出现的苏美尔语名称表明，它们建立在更早但如今已无直接文献留存的早期青铜时代苏美尔传统之上。$^{[30]}$

古典黄道体系是对公元前6世纪新巴比伦星座的修订。希腊人在公元前4世纪采纳了巴比伦的星座体系。托勒密星座中有二十个源自古代近东，另有十个使用相同的恒星，但名称不同。$^{[29]}$

《圣经》学者 E.\,W.\,Bullinger 将《以西结书》和《启示录》中提及的一些生物解释为黄道四分区的中间星座标志$^{[31][32]}$：狮子对应狮子座，公牛对应金牛座，人象征宝瓶座，而鹰则代表天蝎座。$^{[33]}$ 《圣经·约伯记》中也提及若干星座，包括 עיש ‘Ayish（意为“灵柩”）、כסיל chesil（“愚者”）以及 כימה chimah（“堆积”）（约伯记 9:9，38:31–32），在《钦定版圣经》中被译为“大角星、猎户座和昴宿星团”，但其中 ‘Ayish“灵柩” 实际对应的是大熊座。$^{[34]}$ 术语 Mazzaroth מַזָּרוֹת，被译为“冠冕的花环”，在《约伯记》38:32中是一个仅出现一次的词，其含义可能指黄道星座。
\subsubsection{古典古代}
\begin{figure}[ht]
\centering
\includegraphics[width=6cm]{./figures/8bd69f4efca5ef5d.png}
\caption{埃及星图与旬星时钟，出自塞嫩穆特墓穴天花板，约公元前1473年} \label{fig_Coti_5}
\end{figure}
关于古希腊星座的信息所存不多，仅有一些零散证据可见于希腊诗人赫西俄德的《工作与时日》，其中提到了“天体”。$^{[35]}$在希腊化时期，希腊天文学基本上采纳了更早的巴比伦体系，该体系最早由克尼多斯的欧多克索斯于公元前4世纪引入希腊。$^{[citation needed]}$ 欧多克索斯的原著已佚，但其内容以诗体形式保存在公元前3世纪阿拉托斯的作品中。现存关于星座神话起源最为完整的著作，出自一位被称为“伪埃拉托色尼”的希腊化时期作家，以及一位被称为“伪许吉努斯”的早期罗马作家。从古代晚期直到近代早期，西方天文学教学的基础是托勒密于公元2世纪撰写的《天文学大成》。

在托勒密王国时期，埃及本土传统以拟人化形象来表现行星、恒星及各种星座。$^{[36]}$ 这些传统后来与希腊和巴比伦的天文学体系相结合，最终形成了丹德拉黄道，这是目前已知最早完整描绘全部现今熟知黄道星座的实例，同时还包含一些原创的埃及星座、旬星以及行星。$^{[28][37]}$ 这一融合发生的具体时间仍不明确，但多数观点认为其形成于罗马时期，即公元2至4世纪之间。托勒密的《天文学大成》在中世纪时期仍是欧洲及伊斯兰天文学中界定星座的标准依据。$^{[38]}$
\subsubsection{中国古代}
\begin{figure}[ht]
\centering
\includegraphics[width=8cm]{./figures/4fe77cf6fcfe09c6.png}
\caption{采用圆柱投影的中国星图（苏颂）} \label{fig_Coti_6}
\end{figure}
中国古代具有悠久的天象观测传统。$^{[39]}$ 一些并不具体的中国恒星名称，后来被归入二十八宿体系，已在出土于安阳、可追溯至商代中期的甲骨文中发现。中国星座体系是对中国天空最为重要的观测成果之一，其历史可追溯至公元前5世纪。中国的星座体系独立于希腊—罗马体系而发展，尽管通过与古巴比伦天文学的若干相似之处，推测二者之间可能存在较早的相互影响。$^{[40]}$

汉代经典中国天文学的三大学派，被认为源自更早的战国时期天文学家。三大学派的星座体系在三国时期、3世纪的天文学家陈卓手中被整合为一个统一体系。陈卓的原著已佚，但其星座体系的信息保存在唐代文献中，尤以居住在长安的中印天文学家瞿昙悉达所著《九执历》（Chinese: 九執曆；pinyin: Jiǔzhí-lì）为代表。$^{[41]}$ 现存最早的中国星图亦可追溯至这一时期，并作为敦煌文献的一部分保存至今。中国本土天文学在宋代达到繁荣，而在元代以及随后的明代，则逐渐受到中世纪时期其他东方和西方天文学家的影响，参见《开元占经》。$^{[40]}$ 由于这一时期绘制的星图采用了更具科学性的方式，因此被认为更为可靠。$^{[42]}$

宋代一幅著名的星图是《苏州天文图》，该图将中国天空的恒星雕刻在一块石制平面天球图上，基于实际观测精确绘制，并标示了位于金牛座的公元1054年超新星。$^{[42]}$

明代晚期在欧洲天文学的影响下，星图中描绘了更多恒星，但仍保留了传统星座体系。新观测到的恒星被作为补充并入南天的旧有星座之中，这些南天区域并未出现在古代中国天文学家的传统记录中。明代后期，徐光启与德国耶稣会士汤若望（Johann Adam Schall von Bell）进一步改进了这一体系，其成果记录于《崇祯历书》（1628年）。$^{[clarification needed]}$ 传统中国星图基于西方星图的知识，引入了南天23个新星座、共125颗恒星；随着这一改进，中国的天空体系得以与世界天文学接轨。$^{[42][43]}$
\subsection{近代早期天文学}  
从历史角度看，北天与南天星座的起源具有显著差异。大多数北天星座可追溯至古代，其名称主要源自古典希腊神话。$^{[11]}$ 关于这些星座的证据以星图形式保存下来，其中最早的表现形式可能出现在被称为“法尔内塞天球仪”的雕像上，该雕像或许基于希腊天文学家喜帕恰斯的星表。$^{[44]}$ 南天星座则多为近代的创造，有时作为古代星座的替代（例如阿尔戈船座）。一些南天星座最初名称较长，后来被缩短为更易使用的形式，例如 Musca Australis 最终简化为 Musca。$^{[11]}$

一些早期提出的星座从未被普遍接受。不同观测者常以不同方式将恒星组合成星座，而任意划定的星座边界也常导致对天体归属星座的混淆。在天文学家开始精确划定星座边界之前（始于19世纪），星座通常只是天空中界限模糊的区域。$^{[45]}$ 如今，星座边界遵循官方认可的赤经与赤纬分界线，这些界线基于本杰明·古尔德在其星表《阿根廷天区测量》（Uranometria Argentina）中所定义的 B1875.0 春分点。$^{[46]}$

1603年，约翰·拜耳出版的星图集《天图》（Uranometria）将恒星明确归属至各个星座，并通过为每个星座内的恒星分配希腊字母与拉丁字母，正式确立了这一划分体系。这种命名方式如今被称为拜耳命名法。$^{[47]}$ 随后的星图编制进一步促成了当今通行的现代星座体系。

若干近代提出的星座方案并未延续下来。例如，法国天文学家皮埃尔·勒莫尼耶和约瑟夫·拉朗德曾提出一些一度流行、但后来被废弃的星座。北天星座“象限仪座”（Quadrans Muralis）曾沿用至19世纪，其名称后来被用于象限仪座流星雨，但如今该星座已被划分并并入牧夫座与天龙座之中。$^{[48]}$
\subsubsection{南天星座的起源}
\begin{figure}[ht]
\centering
\includegraphics[width=8cm]{./figures/8f18b38336c2466a.png}
\caption{葡萄牙天文学家若昂·法拉斯于1500年5月1日绘制的南天球天空素描} \label{fig_Coti_7}
\end{figure}
\begin{figure}[ht]
\centering
\includegraphics[width=8cm]{./figures/0786e3410b2ca56b.png}
\caption{荷兰制图学黄金时代的天球图，由荷兰制图师弗雷德里克·德·维特绘制} \label{fig_Coti_8}
\end{figure}
赤纬约低于−65°的南天天空，仅被古代北方的巴比伦、埃及、希腊、中国和波斯天文学家部分编录。北天与南天恒星图案存在差异这一认知，可追溯至古典时期作家的记载，例如他们描述了约公元前600年由埃及法老尼科二世下令进行的非洲环航探险，以及约公元前500年航海家汉诺的航行。

南天星座的历史并非一条清晰而单一的发展脉络。不同的观测者提出了不同的分组方式和不同的名称，其中一些反映了各自的民族传统，另一些则旨在彰显或纪念特定的资助者。自14世纪至16世纪，南天星座具有重要意义，因为航海者依赖恒星进行天文导航。记录了新南天星座的意大利探险者包括安德烈亚·科尔萨利、安东尼奥·皮加费塔以及亚美利哥·韦斯普奇。$^{[33]}$

在这一地区，国际天文学联合会承认的88个星座中，有许多最早出现在16世纪末引入的天球仪上，这些天球仪主要由彼得鲁斯·普朗修斯绘制，依据的是荷兰航海家彼得·迪尔克松·凯泽$^{[49]}$和弗雷德里克·德·豪特曼$^{[50][51][52][53]}$的观测结果。这些星座通过约翰·拜耳于1603年出版的星图集《天图》（Uranometria）而广为人知。$^{[54]}$ 此外，法国天文学家尼古拉·路易·德·拉卡伊于1752年又创建了14个星座，并将古老的阿尔戈船座拆分为三个星座；这些新星座出现在他于1756年出版的星表中。$^{[55]}$
\subsubsection{88个现代星座} 
1922年，国际天文学联合会制定了一份包含88个星座的名单。$^{[56]}$ 该名单大体基于托勒密在公元2世纪《天文学大成》中列出的传统希腊星座以及阿拉托斯的《天象志》，并结合了近代早期的修订与补充（其中最重要的是引入了覆盖托勒密未知的南天天区的星座），这些修订主要由彼得鲁斯·普朗修斯（1592年、1597/98年和1613年）、约翰内斯·赫维留斯（1690年）以及尼古拉·路易·德·拉卡伊（1763年）完成。$^{[57][58][59]}$ 拉卡伊共引入了14个新星座。$^{[60]}$ 1751年至1752年间，拉卡伊在好望角对南半球恒星进行了系统研究，据称他使用口径为0.5英寸（13\,mm）的折射望远镜观测了超过一万颗恒星。

1922年，亨利·诺里斯·罗素编制了一份包含88个星座的名单，并为其指定了三字母缩写。$^{[61]}$ 然而，这些星座之间尚未划定明确的边界。1928年，国际天文学联合会正式接受了88个现代星座，并采纳了由欧仁·德尔波特制定的、沿赤经与赤纬的垂直和水平线划分的连续边界，这些边界共同覆盖了整个天球。$^{[5][62][63]}$ 该名单最终于1930年正式公布。$^{[4]}$ 在可能的情况下，这些现代星座通常沿用了其希腊—罗马前身的名称，例如猎户座、狮子座或天蝎座。该体系的目标在于进行面积划分，即将整个天球划分为相互连续的区域。$^{[57]}$ 在这88个现代星座中，有36个主要位于北天，其余52个主要位于南天。
\begin{figure}[ht]
\centering
\includegraphics[width=8cm]{./figures/90ffc7c404c16dde.png}
\caption{88个国际天文学联合会指定星座的等距圆柱投影图，其中黄道星座以轮廓标出} \label{fig_Coti_9}
\end{figure}
德尔波特所制定的星座边界使用了可追溯至 B1875.0 历元的数据，当时本杰明·A·古尔德首次提出为天球划定边界的设想，德尔波特的工作正是基于这一建议展开的。$^{[64]}$由于该历元较早，其直接后果是：受岁差影响，在现代星图（例如 J2000 历元）上，这些边界已经出现一定程度的倾斜，不再完全呈现为严格的垂直或水平。$^{[65]}$随着未来数十年乃至数百年的推移，这一效应还将进一步加剧。
\subsubsection{符号}  
星座本身并没有官方规定的符号，不过位于黄道上的星座可以使用黄道十二宫的符号。$^{[66]}$ 其他现代星座以及一些仍在现代命名体系中出现的古老星座，其符号曾在个别情况下被发表过。$^{[67][68]}$
\subsection{暗云星座}
\begin{figure}[ht]
\centering
\includegraphics[width=8cm]{./figures/b131a70b41d9ca64.png}
\caption{从地球观测到的银河系，其中显著的暗云结构以白色标注，显著的恒星云以黑色标注。图像下方偏左的位置可见大麦哲伦云和小麦哲伦云。呈星状的天体为银河系的球状星团。} \label{fig_Coti_10}
\end{figure}
银河裂隙是银河中一系列暗色斑块，在南天天空中最为显著。$^{[69]}$ 一些文化在这些暗斑中辨识出特定的形状。印加文明的成员将银河中的若干暗区或暗星云视为动物，并将其出现与季节性降雨联系起来。$^{[70][71][72]}$ 澳大利亚原住民天文学同样描述了暗云星座，其中最著名的是“天空中的鸸鹋”，其头部并非由恒星构成，而是由一团暗星云——煤袋星云——所形成。$^{[73]}$
\begin{figure}[ht]
\centering
\includegraphics[width=8cm]{./figures/0c99d553b7af22a3.png}
\caption{“天空中的鸸鹋”——一种由暗云而非恒星定义的星座。鸸鹋的头部是煤袋星云，其正上方为南十字座，左侧为天蝎座。} \label{fig_Coti_11}
\end{figure}
\begin{figure}[ht]
\centering
\includegraphics[width=8cm]{./figures/0952d0820954f62f.png}
\caption{印加文明在 Mayu（天河），即银河中的暗云星座；南十字座位于 Yutu 的上方，而美洲驼的双眼则是半人马座的 $\alpha$ 星与 $\beta$ 星。} \label{fig_Coti_13}
\end{figure}
\subsubsection{暗云星座列表}
\begin{figure}[ht]
\centering
\includegraphics[width=6cm]{./figures/b8843b43bf396e6c.png}
\caption{奥斯曼时期的宇宙志地图，描绘了行星天球、黄道十二宫以及月宿（《Zubdat al-Tawarikh》）} \label{fig_Coti_14}
\end{figure}
\begin{itemize}
\item 银河裂隙（天文学）  
\item 天鹅座裂隙  
\item 巨蛇座—天鹰座裂隙  
\item 暗马（天文学）  
\item 蛇夫座 $\rho$ 云复合体  
\item 天空中的鸸鹋  
\end{itemize}
\subsection{参见}  
\begin{itemize}
\item 天球制图学  
\item 星座家族  
\item 已废弃的星座  
\item 按星座列出的恒星列表  
\item 约翰内斯·赫维留斯所列星座  
\item 拉卡伊所列星座  
\item 彼得鲁斯·普朗修斯所列星座  
引用

 “星座的定义”。牛津英语词典。于2013年1月2日从原文存档。检索日期：2016年8月2日
 汉娜，罗伯特（2021）。“古希腊的星辰”。载于：鲍西卡斯，埃夫罗西尼；麦克拉斯基，斯蒂芬·C.；斯蒂尔，约翰（编）。推进文化天文学：致敬克莱夫·拉格尔斯的研究。瑞士萨姆、楚格：施普林格自然出版社。第212–217. ISBN 978-3-030-64606-6.
 布里顿，约翰·P.（2010）。“巴比伦月球理论研究：第三部分——均匀黄道的引入”。精确科学史档案. 64 (6): 617–63. doi:10.1007/S00407-010-0064-Z. JSTOR 41134332. S2CID 122004678. [T]黄道于公元前408年至公元前397年间引入，很可能是在公元前400年后的短短几年内。
 德尔波尔特，欧仁（1930）。星座的科学界定. 国际天文学联合会.
 里德帕斯 2018, 第4.
 “DOCdb深空观测者伴侣——在线数据库”。检索日期：2018年9月21日。
 “星群完整列表”。于2012年9月29日从原文存档。检索日期：2018年9月21日。
 “星座 | 在线词源词典中星座的起源与含义”. www.etymonline.com.
 “星座”.牛津天文学词典. 检索日期2019年7月26日
 哈博德，约翰·布拉德利；古德温，H. B.（1897）。航海术语表：实用航海者的随身手册（第3版）。朴茨茅斯：格里芬。第142页。
 诺顿，阿瑟·P.（1959）。诺顿星图。第1页。
 斯蒂尔，乔尔·多尔曼（1884）恒星的故事：新式描述天文学, 科学丛书，美国图书公司，第220页
 “星座会分裂或改变吗？”美国国家航空航天局。于2011年10月13日从原文存档。检索日期：2014年11月27日。
 库佩利斯，西奥；库恩，卡尔·F.（2007）。追寻宇宙。琼斯与巴特利特出版社。第369页。ISBN 978-0-7637-4387-1.
 科瓦列夫斯基，让；塞德尔曼，P·肯尼思（2004）。天体测量学基础。剑桥大学出版社。ISBN 978-0-521-64216-3.
 索费尔，M；克利奥纳，S. A；佩蒂，G；沃尔夫，P；科佩金，S. M；布雷塔尼翁，P；布伦贝格，V. A；卡皮坦，N；达穆尔，T；福岛，T；吉诺，B；黄廷宇；林德格伦，L；马晓东；诺德特韦特，K；赖斯，J. C；塞德尔曼，P. K；沃克鲁赫利克，D；威尔，C.M；徐，C（2003）。“国际天文学联合会2000年关于相对论框架下天体测量、天体力学和计量学的决议：说明性补充”。天文期刊. 126 (6): 2687–706. arXiv:astro-ph/0303376. 文献标识码:2003AJ....126.2687S. DOI:10.1086/378162. S2CID 32887246.
 “关于光谱“质心径向速度测量”定义的决议C1”。特刊：2003年7月13日至26日在悉尼举行的第25届国际天文学联合会大会初步议程，信息公告第91号（PDF）。国际天文学联合会秘书处。2002年7月。第50页。存档于2022年10月9日的原始版本。检索日期：2017年9月28日。
 “林中之树——我们为何能识别面孔与星座”.鹦鹉螺，2014年5月19日。检索日期：2020年2月3日
 雷伊，H. A. (1954). 《群星：一种全新的观星方式》。霍顿·米夫林·哈考特出版公司。ISBN 978-0547132808.
 “按季节精选的最佳星座：南半球”.星图指南. 2023年6月22日. 检索日期2024年6月24日
 戴尔，艾伦（2016年8月19日）。“南半球的移动星辰”。奇妙的天空。检索日期：2024年6月24日
 查特兰，马克·R.（1991）。奥杜邦学会夜空野外指南。纽约：A. A. 克诺夫。第134–162、405–420页。ISBN978-0-679-73354-6。
 萨特，保罗·M.（2020年10月9日）。“太阳系存在第二个平面，天体围绕太阳运行”。今日宇宙。检索日期：2024年7月5日
 “一月星座”.星座指南. 检索日期2024年6月24日
 “7月星座”.星座指南. 检索日期2024年6月24日
 拉彭格吕克，M.（1997）。“拉卡斯洞窟‘公牛厅’中的昴星团。拉卡斯洞窟的岩画是否展现了马格德林时期（约公元前15300年）的昴星团开放星团？”". 《天文学与文化》: 217. 文献编码:1997ascu.conf..217R.
 坎宁安，D.（2011）。“世界上最古老的地图：破译西班牙埃尔卡斯蒂略洞穴和法国拉斯科洞穴的手绘图案”。午夜科学. 4: 3.
 罗杰斯，J·H（1998）。“古代星座的起源：I.美索不达米亚传统”。英国天文协会期刊. 108: 9. 文献标识码:1998JBAA..108....9R.
 谢弗，布拉德利·E.（2006）。“希腊星座的起源”。科学美国人. 295 (5): 96–101. 天文学文献标识码:2006SciAm.295e..96S. DOI:10.1038/scientificamerican1106-96. PMID 17076089.
 “星座与星名的历史——第4部分：苏美尔的星座与星名？”.加里·D·汤普森。2015年4月21日。于2015年9月7日从原文存档。检索日期：2015年8月30日
 E·威廉·布林格（2015）。星辰的见证。电子书项目。ISBN 978-963-527-403-1.
 丹尼斯·詹姆斯·肯尼迪（1989年6月）。黄道十二宫的真正含义。珊瑚岭事工媒体公司ISBN 978-1-929626-14-4.
 理查德·H·艾伦（2013）。恒星名称：它们的传说与含义。柯瑞尔公司。ISBN 978-0-486-13766-7.
 “H5906 - ʿayiš - 斯特朗希伯来语词典（钦定版）”. 蓝字圣经.
 洛里默，H. L.（1951）。“荷马与赫西俄德笔下的恒星与星座”。雅典英国学校年鉴。46：86–101。doi10.1017/S0068245400018359S2CID192976174
 马歇尔·克拉吉特（1989）。古埃及科学：日历、钟表与天文学。美国哲学学会。ISBN 978-0-87169-214-6.
 丹德拉（1825年）。丹德拉黄道十二宫，概要。（展览，莱斯特广场）.
 伊恩·里德帕斯。“星象故事——托勒密”。检索日期：2025年11月3日。
 尼德姆，约瑟夫（1959）。数学与天文学和地球科学。中国科学与文明。第3卷。剑桥大学出版社。第171页。ISBN 978-0521058018.
 孙晓春；雅各布·基斯特梅克（1997）。汉代的中国天空：星象与社会。布里尔出版社。ISBN 978-90-04-10737-3.
 马克，比尔·M.（2023）。“公元8世纪印度天文学论著的中文译本：瞿昙西达所著《九执经》”。传统欧亚学术中的多语现象。布里尔出版社。第352页。DOI：10.1163/9789004527256_031。ISBN978-90-04-52725-6
 塞林，海莱恩·伊丽丝（2008）。非西方文化中的科学、技术与医学史百科全书。施普林格科学与商业媒体。第2022页。ISBN 978-1-4020-4559-2.
 孙晓春（1997）。海琳·塞林（编）。非西方文化中的科学、技术与医学史百科全书。克鲁尔学术出版社。第910页。ISBN 978-0-7923-4066-9.
 谢弗，布拉德利·E.（2005年5月）。“法尔内塞星图上星座的纪元及其在喜帕恰斯失传目录中的起源” （PDF）. 天文学史杂志. 36/2 (123): 167–19. 文献标识码:2005JHA....36..167S. DOI:10.1177/002182860503600202. S2CID 15431718. 已归档 （PDF）自原始网页于2022年10月9日。
 诺顿，阿瑟·P.（1919）。诺顿星图。第1页。
 “天文纪元”。自原始版本存档于2011年7月24日。检索日期：2010年7月16日。
 斯韦德洛，N. M.（1986年8月）：“约翰内斯·拜耳所用的星表”。天文学史杂志. 17 (5): 189–97. 文献标识码:1986JHA....17..189S. DOI:10.1177/002182868601700304. S2CID 118829690.
 巴伦廷，约翰（2016）。失落的星座。纽约：施普林格。第304页。ISBN 978-3-319-22794-8.
 霍格，海伦·索耶（1951）。“从旧书中寻觅（彼得·迪尔克兹·凯泽，南方星座的描绘者）”。加拿大皇家天文学会会刊. 45: 215. 文献标识码:1951JRASC..45..215S.
 诺贝尔，E. B.（1917）。“论弗雷德里克·德·豪特曼的南天星表及南天星座的起源”，皇家天文学会月报, 77, 414–32 ADS链接.
 德克尔，埃莉（1987）：“南天早期探索”，科学年鉴, 44, 439–70 DOI链接.
 德克尔，埃莉（1987）：“论南天星空知识的传播”，地球之友, 35–37, 211–30 JSTOR链接.
 韦本特，弗兰克；范亨特，罗伯特·H.（2011）。“南天早期星表：德·豪特曼、开普勒（第二类和第三类）以及哈雷”，天文学与天体物理学, 530, A93ADS链接.
 伊恩·里德帕斯.“约翰·拜耳的南天星图”. 星空故事.
 伊恩·里德帕斯。“拉凯耶1756年南天星图”. 星空传说.
 里德帕斯，伊恩（2018）。“星空传说：最后的88个星座”.
 “星座”。国际天文学联合会——IAU。检索日期：2015年8月29日。
 伊恩·里德帕斯。“星座名称、缩写及大小”。检索日期：2015年8月30日。
 伊恩·里德帕斯。“星图——托勒密的《天球图》”。检索日期：2015年8月30日。
 伊恩·里德帕斯.“拉凯伊在开普敦”。检索日期：2022年7月4日。
 “1922年星座的原始名称与缩写”. 检索日期2010年1月31日.
 “星座边界”. 检索日期2011年5月24日.
 拉希耶-雷，马克；吕米内，让-皮埃尔（2001）。天体宝库：从天球之音到征服太空。法国巴黎：剑桥大学出版社，法国国家图书馆。第80页。ISBN 978-0-521-80040-2.
 伊恩·里德帕斯。“本杰明·阿普瑟普·古尔德与《阿根廷星图》”。星空传说
 A.C. 戴文霍尔与S.K. 莱格特，“星座边界数据目录”，（斯特拉斯堡天文数据中心，1990年2月）。
 例如，在1833年航海历与天文年历（伦敦海军部）中
 彼得·格雷戈（2012）星空之书：北半球四季观星指南。F+W传媒。
 米勒，柯克（2024年10月18日）。“星座符号的初步展示”（PDF）。unicode.org。Unicode联盟。检索日期：2024年10月22日
 饶，乔（2009年9月11日）。“观赏银河的绝佳一周”。太空。检索日期：2024年7月22日。
 “夜空”.Astronomy.pomona.edu。于2010年12月16日从原始网页存档。检索日期：2019年3月12日
 迪尔伯恩，D.S.P.；怀特，R.E.（1983）。“马丘比丘的‘托雷翁’作为天文台”。考古天文学.14(5)：S37。Bibcode:1983JHAS...14...37D
 克鲁普，埃德温（1994）。古代天空的回响。米尼奥拉：多佛出版公司，第47至51页。ISBN978-0486428826。
 博尔德洛，安德烈·G.（2013）。夜空的旗帜：当天文学邂逅民族自豪感。施普林格科学与商业媒体。第124页及以后。ISBN 978-1-4614-0929-8.托勒密所列星座  
\end{itemize}
\subsection{参考文献} 
\textbf{注释}

a.占星学引用

\textbf{引用}
\begin{enumerate}
“星座的定义”。牛津英语词典。于2013年1月2日从原文存档。检索日期：2016年8月2日
 汉娜，罗伯特（2021）。“古希腊的星辰”。载于：鲍西卡斯，埃夫罗西尼；麦克拉斯基，斯蒂芬·C.；斯蒂尔，约翰（编）。推进文化天文学：致敬克莱夫·拉格尔斯的研究。瑞士萨姆、楚格：施普林格自然出版社。第212–217. ISBN 978-3-030-64606-6.
 布里顿，约翰·P.（2010）。“巴比伦月球理论研究：第三部分——均匀黄道的引入”。精确科学史档案. 64 (6): 617–63. doi:10.1007/S00407-010-0064-Z. JSTOR 41134332. S2CID 122004678. [T]黄道于公元前408年至公元前397年间引入，很可能是在公元前400年后的短短几年内。
 德尔波尔特，欧仁（1930）。星座的科学界定. 国际天文学联合会.
 里德帕斯 2018, 第4.
 “DOCdb深空观测者伴侣——在线数据库”。检索日期：2018年9月21日。
 “星群完整列表”。于2012年9月29日从原文存档。检索日期：2018年9月21日。
 “星座 | 在线词源词典中星座的起源与含义”. www.etymonline.com.
 “星座”.牛津天文学词典. 检索日期2019年7月26日
 哈博德，约翰·布拉德利；古德温，H. B.（1897）。航海术语表：实用航海者的随身手册（第3版）。朴茨茅斯：格里芬。第142页。
 诺顿，阿瑟·P.（1959）。诺顿星图。第1页。
 斯蒂尔，乔尔·多尔曼（1884）恒星的故事：新式描述天文学, 科学丛书，美国图书公司，第220页
 “星座会分裂或改变吗？”美国国家航空航天局。于2011年10月13日从原文存档。检索日期：2014年11月27日。
 库佩利斯，西奥；库恩，卡尔·F.（2007）。追寻宇宙。琼斯与巴特利特出版社。第369页。ISBN 978-0-7637-4387-1.
 科瓦列夫斯基，让；塞德尔曼，P·肯尼思（2004）。天体测量学基础。剑桥大学出版社。ISBN 978-0-521-64216-3.
 索费尔，M；克利奥纳，S. A；佩蒂，G；沃尔夫，P；科佩金，S. M；布雷塔尼翁，P；布伦贝格，V. A；卡皮坦，N；达穆尔，T；福岛，T；吉诺，B；黄廷宇；林德格伦，L；马晓东；诺德特韦特，K；赖斯，J. C；塞德尔曼，P. K；沃克鲁赫利克，D；威尔，C.M；徐，C（2003）。“国际天文学联合会2000年关于相对论框架下天体测量、天体力学和计量学的决议：说明性补充”。天文期刊. 126 (6): 2687–706. arXiv:astro-ph/0303376. 文献标识码:2003AJ....126.2687S. DOI:10.1086/378162. S2CID 32887246.
 “关于光谱“质心径向速度测量”定义的决议C1”。特刊：2003年7月13日至26日在悉尼举行的第25届国际天文学联合会大会初步议程，信息公告第91号（PDF）。国际天文学联合会秘书处。2002年7月。第50页。存档于2022年10月9日的原始版本。检索日期：2017年9月28日。
 “林中之树——我们为何能识别面孔与星座”.鹦鹉螺，2014年5月19日。检索日期：2020年2月3日
 雷伊，H. A. (1954). 《群星：一种全新的观星方式》。霍顿·米夫林·哈考特出版公司。ISBN 978-0547132808.
 “按季节精选的最佳星座：南半球”.星图指南. 2023年6月22日. 检索日期2024年6月24日
 戴尔，艾伦（2016年8月19日）。“南半球的移动星辰”。奇妙的天空。检索日期：2024年6月24日
 查特兰，马克·R.（1991）。奥杜邦学会夜空野外指南。纽约：A. A. 克诺夫。第134–162、405–420页。ISBN978-0-679-73354-6。
 萨特，保罗·M.（2020年10月9日）。“太阳系存在第二个平面，天体围绕太阳运行”。今日宇宙。检索日期：2024年7月5日
 “一月星座”.星座指南. 检索日期2024年6月24日
 “7月星座”.星座指南. 检索日期2024年6月24日
 拉彭格吕克，M.（1997）。“拉卡斯洞窟‘公牛厅’中的昴星团。拉卡斯洞窟的岩画是否展现了马格德林时期（约公元前15300年）的昴星团开放星团？”". 《天文学与文化》: 217. 文献编码:1997ascu.conf..217R.
 坎宁安，D.（2011）。“世界上最古老的地图：破译西班牙埃尔卡斯蒂略洞穴和法国拉斯科洞穴的手绘图案”。午夜科学. 4: 3.
 罗杰斯，J·H（1998）。“古代星座的起源：I.美索不达米亚传统”。英国天文协会期刊. 108: 9. 文献标识码:1998JBAA..108....9R.
 谢弗，布拉德利·E.（2006）。“希腊星座的起源”。科学美国人. 295 (5): 96–101. 天文学文献标识码:2006SciAm.295e..96S. DOI:10.1038/scientificamerican1106-96. PMID 17076089.
 “星座与星名的历史——第4部分：苏美尔的星座与星名？”.加里·D·汤普森。2015年4月21日。于2015年9月7日从原文存档。检索日期：2015年8月30日
 E·威廉·布林格（2015）。星辰的见证。电子书项目。ISBN 978-963-527-403-1.
 丹尼斯·詹姆斯·肯尼迪（1989年6月）。黄道十二宫的真正含义。珊瑚岭事工媒体公司ISBN 978-1-929626-14-4.
 理查德·H·艾伦（2013）。恒星名称：它们的传说与含义。柯瑞尔公司。ISBN 978-0-486-13766-7.
 “H5906 - ʿayiš - 斯特朗希伯来语词典（钦定版）”. 蓝字圣经.
 洛里默，H. L.（1951）。“荷马与赫西俄德笔下的恒星与星座”。雅典英国学校年鉴。46：86–101。doi10.1017/S0068245400018359S2CID192976174
 马歇尔·克拉吉特（1989）。古埃及科学：日历、钟表与天文学。美国哲学学会。ISBN 978-0-87169-214-6.
 丹德拉（1825年）。丹德拉黄道十二宫，概要。（展览，莱斯特广场）.
 伊恩·里德帕斯。“星象故事——托勒密”。检索日期：2025年11月3日。
 尼德姆，约瑟夫（1959）。数学与天文学和地球科学。中国科学与文明。第3卷。剑桥大学出版社。第171页。ISBN 978-0521058018.
 孙晓春；雅各布·基斯特梅克（1997）。汉代的中国天空：星象与社会。布里尔出版社。ISBN 978-90-04-10737-3.
 马克，比尔·M.（2023）。“公元8世纪印度天文学论著的中文译本：瞿昙西达所著《九执经》”。传统欧亚学术中的多语现象。布里尔出版社。第352页。DOI：10.1163/9789004527256_031。ISBN978-90-04-52725-6
 塞林，海莱恩·伊丽丝（2008）。非西方文化中的科学、技术与医学史百科全书。施普林格科学与商业媒体。第2022页。ISBN 978-1-4020-4559-2.
 孙晓春（1997）。海琳·塞林（编）。非西方文化中的科学、技术与医学史百科全书。克鲁尔学术出版社。第910页。ISBN 978-0-7923-4066-9.
 谢弗，布拉德利·E.（2005年5月）。“法尔内塞星图上星座的纪元及其在喜帕恰斯失传目录中的起源” （PDF）. 天文学史杂志. 36/2 (123): 167–19. 文献标识码:2005JHA....36..167S. DOI:10.1177/002182860503600202. S2CID 15431718. 已归档 （PDF）自原始网页于2022年10月9日。
 诺顿，阿瑟·P.（1919）。诺顿星图。第1页。
 “天文纪元”。自原始版本存档于2011年7月24日。检索日期：2010年7月16日。
 斯韦德洛，N. M.（1986年8月）：“约翰内斯·拜耳所用的星表”。天文学史杂志. 17 (5): 189–97. 文献标识码:1986JHA....17..189S. DOI:10.1177/002182868601700304. S2CID 118829690.
 巴伦廷，约翰（2016）。失落的星座。纽约：施普林格。第304页。ISBN 978-3-319-22794-8.
 霍格，海伦·索耶（1951）。“从旧书中寻觅（彼得·迪尔克兹·凯泽，南方星座的描绘者）”。加拿大皇家天文学会会刊. 45: 215. 文献标识码:1951JRASC..45..215S.
 诺贝尔，E. B.（1917）。“论弗雷德里克·德·豪特曼的南天星表及南天星座的起源”，皇家天文学会月报, 77, 414–32 ADS链接.
 德克尔，埃莉（1987）：“南天早期探索”，科学年鉴, 44, 439–70 DOI链接.
 德克尔，埃莉（1987）：“论南天星空知识的传播”，地球之友, 35–37, 211–30 JSTOR链接.
 韦本特，弗兰克；范亨特，罗伯特·H.（2011）。“南天早期星表：德·豪特曼、开普勒（第二类和第三类）以及哈雷”，天文学与天体物理学, 530, A93ADS链接.
 伊恩·里德帕斯.“约翰·拜耳的南天星图”. 星空故事.
 伊恩·里德帕斯。“拉凯耶1756年南天星图”. 星空传说.
 里德帕斯，伊恩（2018）。“星空传说：最后的88个星座”.
 “星座”。国际天文学联合会——IAU。检索日期：2015年8月29日。
 伊恩·里德帕斯。“星座名称、缩写及大小”。检索日期：2015年8月30日。
 伊恩·里德帕斯。“星图——托勒密的《天球图》”。检索日期：2015年8月30日。
 伊恩·里德帕斯.“拉凯伊在开普敦”。检索日期：2022年7月4日。
 “1922年星座的原始名称与缩写”. 检索日期2010年1月31日.
 “星座边界”. 检索日期2011年5月24日.
 拉希耶-雷，马克；吕米内，让-皮埃尔（2001）。天体宝库：从天球之音到征服太空。法国巴黎：剑桥大学出版社，法国国家图书馆。第80页。ISBN 978-0-521-80040-2.
 伊恩·里德帕斯。“本杰明·阿普瑟普·古尔德与《阿根廷星图》”。星空传说
 A.C. 戴文霍尔与S.K. 莱格特，“星座边界数据目录”，（斯特拉斯堡天文数据中心，1990年2月）。
 例如，在1833年航海历与天文年历（伦敦海军部）中
 彼得·格雷戈（2012）星空之书：北半球四季观星指南。F+W传媒。
 米勒，柯克（2024年10月18日）。“星座符号的初步展示”（PDF）。unicode.org。Unicode联盟。检索日期：2024年10月22日
 饶，乔（2009年9月11日）。“观赏银河的绝佳一周”。太空。检索日期：2024年7月22日。
 “夜空”.Astronomy.pomona.edu。于2010年12月16日从原始网页存档。检索日期：2019年3月12日
 迪尔伯恩，D.S.P.；怀特，R.E.（1983）。“马丘比丘的‘托雷翁’作为天文台”。考古天文学.14(5)：S37。Bibcode:1983JHAS...14...37D
 克鲁普，埃德温（1994）。古代天空的回响。米尼奥拉：多佛出版公司，第47至51页。ISBN978-0486428826。
 博尔德洛，安德烈·G.（2013）。夜空的旗帜：当天文学邂逅民族自豪感。施普林格科学与商业媒体。第124页及以后。ISBN 978-1-4614-0929-8.
\end{enumerate}